\subsection{The Daniell extension}

The content of this subsection follows \cite{Yamagami}. As that book states, the
facts that $L^1$ is a Riesz space and that $I$ extends to $L^1$ follow the
procedure of \cite{BM}.
\textbf{Nothing in this subsection is new; it is here to fix the definitions,
notation and assertions used in later subsections.}
Accordingly we state only definitions and assertions below, and refer to the
literature above for proofs.

\subsubsection{Definition of the Daniell extension}

\begin{definition}\label{def:series-representation}
  Let $(X,L,I)$ be an integration system on a set $X$.
  Given a function $f:X\to\mathbb{R}$, if there exist sequences of functions
  $\{f_n\}\subset L$ and $\{\varphi_n\}\subset L$ satisfying the following, then
  the pair of sequences $\{f_n\},\{\varphi_n\}$ is called a \textbf{series
  representation} of the function $f$, and a function admitting a series
  representation is said to be \textbf{integrable} with respect to the
  integration system $(X,L,I)$.
  \begin{enumerate}
    \item $|f_n|\le\varphi_n$ for every $n$;
    \item $\sum_{n=1}^\infty I(\varphi_n)<+\infty$;
    \item for every $x\in X$, if $\sum_{n=1}^\infty \varphi_n(x)<+\infty$ then
      $f(x)=\sum_{n=1}^\infty f_n(x)$.
  \end{enumerate}
  We write a series representation $\{f_n\},\{\varphi_n\}$ of a function $f$ as
  \[
    f\overset{\varphi_n}{\simeq}\sum f_n .
  \]
\end{definition}

\begin{definition}
  Let $(X,L,I)$ be an integration system on a set $X$.
  We denote the set of all functions integrable with respect to the integration
  system $(X,L,I)$ by
  \[
    L^1:=\{f\in\mathbb{R}^X \mid\text{$f$ is integrable with respect to the integration system $(X,L,I)$}\}
  \]
  \footnote{
    $L^1$ depends on $I$ as well, but in this paper we never put several
    integrals on one Riesz space $L$.
    We therefore write $L^1$, as something determined by $L$.
    Where this could be ambiguous we write it out, as in $(L^{\mathrm{M}})^1$.
    When referring to the object itself we use the functorial notation
    $\mathcal{D}(X,L,I)=(X,L^1,I^1)$.
  }.
\end{definition}

\medskip
Definition~\ref{def:series-representation} takes the form in which the
dominating functions $\varphi_n$ are made explicit, but this is equivalent to the
form without dominating functions adopted in \cite{BM}.
We will not use this below, but we state it in order to make the correspondence
with the literature clear.

\begin{lemma}[equivalence of the definitions with and without dominating functions]\label{lem:series-rep-equiv}
  Let $(X,L,I)$ be an integration system on a set $X$ and let
  $f:X\to\mathbb{R}$.
  Then the following are equivalent.
  \begin{enumerate}
    \item[(i)] $f$ is integrable in the sense of
      Definition~\ref{def:series-representation}.
      That is, there exist $\{f_n\},\{\varphi_n\}\subset L$ with
      $|f_n|\le\varphi_n$ and $\sum_n I(\varphi_n)<+\infty$ such that
      $f(x)=\sum_n f_n(x)$ holds at every $x$ with
      $\sum_n\varphi_n(x)<+\infty$.
    \item[(ii)] There exists $\{g_m\}\subset L$ with
      $\sum_m I(|g_m|)<+\infty$ such that $f(x)=\sum_m g_m(x)$ holds at every
      $x$ with $\sum_m|g_m(x)|<+\infty$.
  \end{enumerate}
\end{lemma}
\begin{proof}
  For $(ii)\Rightarrow(i)$ take $f_n:=g_n$ and $\varphi_n:=|g_n|$.
  Since $L$ is a Riesz space we have $|g_n|\in L$, and
  $\sum_n\varphi_n(x)<+\infty$ is literally $\sum_n|g_n(x)|<+\infty$.

  We show $(i)\Rightarrow(ii)$. Put $\psi_n:=\varphi_n-|f_n|\ge0$, so that
  $\psi_n\in L$.
  Define the sequence $\{g_m\}$ by
  \[
    f_1,\ \psi_1,\ -\psi_1,\ f_2,\ \psi_2,\ -\psi_2,\ f_3,\ \psi_3,\ -\psi_3,\ \dots .
  \]
  First, $I(|f_n|)\ge0$ because $I$ preserves order, and
  \[
    \sum_m I(|g_m|)
    =\sum_n\bigl(I(|f_n|)+2I(\psi_n)\bigr)
    =\sum_n\bigl(2I(\varphi_n)-I(|f_n|)\bigr)
    \le2\sum_n I(\varphi_n)<+\infty.
  \]
  Next, if $x$ satisfies $\sum_m|g_m(x)|<+\infty$, then
  \[
    \sum_m|g_m(x)|
    =\sum_n\bigl(|f_n(x)|+2\psi_n(x)\bigr)
    \ge\sum_n\bigl(|f_n(x)|+\psi_n(x)\bigr)
    =\sum_n\varphi_n(x)
  \]
  gives $\sum_n\varphi_n(x)<+\infty$. Hence $f(x)=\sum_n f_n(x)$ by condition
  (i).
  On the other hand $\sum_m g_m(x)$ converges absolutely, so rearrangement and
  regrouping are free, and $\psi_n$ cancels against $-\psi_n$, giving
  \[
    \sum_m g_m(x)=\sum_n f_n(x)=f(x) .
  \]
  Therefore (ii) holds.
\end{proof}

\subsubsection{Construction of the extended integral}

\begin{lemma}\label{lem:equiv-condition-for-riesz-space}
  Let $X$ be a set and let $L\subset\mathbb{R}^X$ be an
  $\mathbb{R}$-vector space.
  Then the condition that $|f|\in L$ for every $f\in L$ is equivalent to $L$
  being a Riesz space on $X$.
\end{lemma}

\begin{lemma}\label{lem:series-exp-is-abs-conv}
  Let $(X,L,I)$ be an integration system on a set $X$ and let
  $f\overset{\varphi_n}{\simeq}\sum f_n$.
  Then, for every $x\in X$, if $\sum_{n=1}^\infty \varphi_n(x)<+\infty$ then
  $f(x)=\sum_{n=1}^\infty f_n(x)$ converges absolutely.
\end{lemma}

\begin{theorem}\label{thm:l1-is-riesz-space}
  Let $(X,L,I)$ be an integration system on a set $X$.
  Then the family $L^1$ of functions integrable with respect to the integration
  system $(X,L,I)$ is a Riesz space on $X$.
\end{theorem}

\begin{lemma}\label{lem:integration-preserve-order}
  Let $(X,L,I)$ be an integration system on a set $X$.
  Then $I(f)\le I(g)$ for all $f,g\in L$ with $f\le g$.
\end{lemma}

\begin{lemma}\label{lem:triangle-ineq-for-integration}
  Let $(X,L,I)$ be an integration system on a set $X$.
  Then $|I(f)|\le I(|f|)$ for every $f\in L$.
\end{lemma}

\begin{lemma}\label{lem:abs-conv-for-int-of-series-rep}
  Let $(X,L,I)$ be an integration system on a set $X$.
  For every integrable function $f\in L^1$ with respect to the integration
  system $(X,L,I)$ and every series representation
  $f\overset{\varphi_n}{\simeq}\sum f_n$ of it, the sum
  $\sum_{n=1}^\infty I(f_n)$ converges absolutely.
\end{lemma}

\begin{lemma}[continuity of the integral]\label{lem:integration-is-continuous}
  Let $(X,L,I)$ be an integration system on a set $X$, let $\{f_n\}$ be
  a sequence of functions in $L$, and let $f\in L$.
  Then, if $f_n(x)\uparrow f(x)$ for every $x\in X$, we have
  $I(f_n)\uparrow I(f)$.
\end{lemma}

\begin{lemma}\label{lem:ext-int-of-zero-func}
  Let $(X,L,I)$ be an integration system on a set $X$.
  Then, given a series representation $0\overset{\eta_n}{\simeq}\sum h_n$ of the
  function constantly equal to $0$,
  \[
    \sum_{n=1}^\infty I(h_n)=0 .
  \]
\end{lemma}

\begin{lemma}\label{lem:ext-int-well-defined}
  Let $(X,L,I)$ be an integration system on a set $X$.
  Then, for any two series representations
  \begin{align*}
    f&\overset{\varphi_n}{\simeq}\sum f_n \\
    f&\overset{\psi_n}{\simeq}\sum g_n
  \end{align*}
  of an integrable function $f\in L^1$ with respect to the integration system
  $(X,L,I)$, we have $\sum_{n=1}^\infty I(f_n)=\sum_{n=1}^\infty I(g_n)$.
\end{lemma}

\begin{definition}
  Let $(X,L,I)$ be an integration system on a set $X$.
  For each integrable function $f\in L^1$ with respect to the integration system
  $(X,L,I)$, choose a series representation
  $f\overset{\varphi_n}{\simeq}\sum f_n$.
  Then the absolutely convergent sum
  (Lemma~\ref{lem:abs-conv-for-int-of-series-rep})
  \[
    \sum_{n=1}^\infty I(f_n)
  \]
  does not depend on the choice of series representation, by
  Lemma~\ref{lem:ext-int-well-defined}.
  We define this value to be $I^1(f)$.
\end{definition}

\begin{lemma}\label{lem:ext-int-is-linear}
  Let $(X,L,I)$ be an integration system on a set $X$.
  Then the map $I^1:L^1\to\mathbb{R}$ is an $\mathbb{R}$-linear map on the Riesz
  space $L^1$ on $X$.
\end{lemma}

\begin{lemma}[truncation estimate for a series representation of a positive function]\label{lem:positive-series-representation-truncation}
  Let $(X,L,I)$ be an integration system on a set $X$ and let
  $f\in L^1$ be an integrable function with $f\ge0$.
  Choose a series representation $f\overset{\varphi_n}{\simeq}\sum f_n$ of $f$
  and, for each $N\ge1$, put
  \[
    S_N:=\sum_{k=1}^Nf_k,\qquad h_N:=S_N\vee0,\qquad q_N:=S_N\wedge0 .
  \]
  Then
  \[
    I(|q_N|)\le\sum_{k=N+1}^\infty I(\varphi_k),
    \qquad
    |I^1(f)-I(h_N)|\le2\sum_{k=N+1}^\infty I(\varphi_k)
  \]
  hold. In particular $I(h_N)\longrightarrow I^1(f)$ as $N\to\infty$.
\end{lemma}

\begin{lemma}\label{lem:ext-int-is-positive}
  Let $(X,L,I)$ be an integration system on a set $X$ and let
  $f\in L^1$.
  Then $f\ge0$ implies $I^1(f)\ge0$.
\end{lemma}

\begin{lemma}\label{lem:positive-good-representation}
  Let $(X,L,I)$ be an integration system on a set $X$ and let
  $f\in L^1$.
  Then, if $f\ge0$, for every $\varepsilon>0$ there is a series representation
  $f\overset{\varphi_n}{\simeq}\sum f_n$ of $f$ satisfying
  \[
    \sum_{n=1}^\infty I(\varphi_n)\le I^1(f)+\varepsilon .
  \]
\end{lemma}

\begin{lemma}[flattening series representations into one sequence]\label{lem:flatten-series-representations}
  Let $(X,L,I)$ be an integration system on a set $X$.
  Suppose that for each $n\ge1$ a series representation
  $u_n\overset{\varphi_{n,k}}{\simeq}\sum_{k=1}^\infty g_{n,k}$ of $u_n\in L^1$
  is given, and that
  \[
    \sum_{n=1}^\infty\sum_{k=1}^\infty I(\varphi_{n,k})<+\infty .
  \]
  Put
  \[
    A:=\left\{x\in X\mathrel{}\middle|\mathrel{}\sum_{n=1}^\infty\sum_{k=1}^\infty\varphi_{n,k}(x)<+\infty\right\} .
  \]
  Then $\sum_{n=1}^\infty u_n(x)$ converges absolutely for every $x\in A$.

  Suppose further that a function $u:X\to\mathbb{R}$ satisfies
  $u(x)=\sum_{n=1}^\infty u_n(x)$ for $x\in A$.
  Choose a bijection $\sigma:\mathbb{N}\to\mathbb{N}\times\mathbb{N}$,
  $\sigma(m)=(n_m,k_m)$, and put $\widetilde{g}_m:=g_{n_m,k_m}$ and
  $\widetilde{\varphi}_m:=\varphi_{n_m,k_m}$.
  Then $u\overset{\widetilde{\varphi}_m}{\simeq}\sum_{m=1}^\infty\widetilde{g}_m$
  is a series representation of $u$.
  In particular $u\in L^1$ and
  \[
    I^1(u)=\sum_{n=1}^\infty I^1(u_n)
  \]
  holds, the right-hand side converging absolutely.
\end{lemma}

\begin{theorem}\label{thm:ext-int-is-integral}
  Let $(X,L,I)$ be an integration system on a set $X$.
  Then the map $I^1:L^1\to\mathbb{R}$ is an integral on the Riesz space $L^1$ on
  $X$.
\end{theorem}

\begin{lemma}\label{lem:daniell-ext-contains-source}
  Let $(X,L,I)$ be an integration system on a set $X$ and let $f\in L$.
  Then $f\in L^1$ and $I^1(f)=I(f)$.
\end{lemma}

\begin{definition}
  Let $(X,L,I)$ be an integration system on a set $X$.
  The integration system $(X,L^1,I^1)$ is called the \textbf{Daniell extension}
  of the integration system $(X,L,I)$.
\end{definition}

\begin{example}[the Lebesgue integral]\label{ex:lebesgue}
  Consider the integration system of Example~\ref{ex:step-functions}, that is,
  the integration system $(X,L,I)$ consisting of the step functions on $[0,1)$
  and the Riemann integral.
  Then its Daniell extension $(X,L^1,I^1)$ coincides with the set of all
  Lebesgue integrable functions on $[0,1)$ together with the Lebesgue integral.

  That is, the construction of Section~\ref{chap:preliminaries} recovers the
  Lebesgue integral.
  This is the origin of the definition of integrability by series
  representations; for details see \cite{Mikusinski} and \cite{BM}, \S2.
\end{example}

\begin{corollary}
  Let $(X,L,I)$ be an integration system on a set $X$ and let
  $(X,L^1,I^1)$ be its Daniell extension.
  Then the identity map $\operatorname{id}_X$ induces a morphism of integration
  systems $i_X:(X,L,I)\to(X,L^1,I^1)$.
\end{corollary}

\subsubsection{The Daniell extension satisfies the MCT}

\begin{lemma}[termwise integration of a nonnegative series]\label{lem:positive-series-flattening}
  Let $(X,L,I)$ be an integration system on a set $X$ and let
  $(X,L^1,I^1)$ be its Daniell extension.
  Suppose $\{h_k\}_{k=1}^{\infty}\subset L^1$ with $h_k\ge0$ satisfies
  $\sum_k I^1(h_k)<\infty$.
  Then, putting
  \[
    h(x):=\begin{cases}
      \sum_n h_n(x) & \text{if $\sum_n h_n(x)$ converges} \\
      0 & \text{otherwise,}
    \end{cases}
  \]
  we have $h\in L^1$ and $I^1(h)=\sum_k I^1(h_k)$.
\end{lemma}

\begin{lemma}\label{lem:daniell-ext-satisfies-mct}
  Let $(X,L,I)$ be an integration system on a set $X$.
  Then its Daniell extension $(X,L^1,I^1)$ satisfies the monotone convergence
  theorem (Definition~\ref{def:satisfies-mct}).
\end{lemma}
